\documentclass{article}
\usepackage{hyperref}
\usepackage[a4paper]{geometry}
\usepackage{fancyhdr}
\pagestyle{fancy}
\lhead{Quantenteilchen}
\rhead{März 2026}
\begin{document}
\section{Quantenteilchen}
Quantenteilchen sind darin besonders, dass die \hyperref[Interferenzmuster mit Quanten]{Interferenzmuster mit Quanten} aufzeigen, dass sie sich sowohl wie Teilchen und wie Wellen, auch mit \hyperref[Die de-Broglie-Wellenlänge]{einer Wellenlänge}, verhalten können \hyperref[Das Delayed-Choice-Experiment]{keinen Ort haben}.
 
\subsection{Begriffe}
Der \emph{Zustand} eines Quantenteilchens beschreibt die bestimmten messbaren Eigenschaften eines Teilchens. Bei Licht, also bei Photonen, ist dies zum Beispiel die Polarisation.
 
Wird durch einen bestimmten Aufbau der Zustand eines Quantenteilchens verändert, ist dies die \emph{Präparation}. Zum Beispiel ein Polarisator, welcher der Polarisation eines Photons festgelegt.
 
Bis es gemessen wird kann ein Quantenteilchen mehrere Zustände gleichzeitig darstellen. Dies ist die \emph{Superposition}.
  
Werden Photonen mit einen Polarisator so präpariert, dass alle die gleiche Polarisation als Zustand haben und treffen dann auf einen Polarisator, welche im $45^\circ$ Winkel zum ersten steht, so haben die Photonen, bis sie auf den diesen auftreffen, sowohl den \textit{Zustand des durchkommens} als auch \textit{den Zustand des nichtdurchkommens} als Superposition.
 
\end{document}